% Shared section of the physics documents: equation of state and 500 nm extinction (\input by each mode document).
\section{The equation of state and the opacity}\label{sec:eos}
Every force balance in this document involves the density $\rho$, and every optical depth involves the extinction
coefficient $\chi_{500}$. Neither is an unknown of its own: at given temperature and pressure both are fixed by the
microphysics of the gas. This section explains how.

\subsection{The ideal gas}
In the solar photosphere and chromosphere the particles interact only weakly, so the gas obeys the ideal-gas law. If $n$
is the number of free particles (atoms, ions, molecules and electrons) per cubic metre,
\begin{equation}
P=n\,k_BT .\label{eq:ideal}
\end{equation}
If the average mass of a free particle is $\bar\mu\,m_{\rm u}$ (with $m_{\rm u}$ the atomic mass unit and $\bar\mu$ the
\textbf{mean molecular weight}), the mass density is $\rho=n\,\bar\mu\,m_{\rm u}$, and therefore
\begin{equation}
\rho=\frac{\bar\mu\,m_{\rm u}}{k_B}\,\frac{P}{T}.\label{eq:rho_ideal}
\end{equation}
At fixed $\bar\mu$ and $T$ the density is proportional to the pressure.

\subsection{\texorpdfstring{Why $\bar\mu$ depends on $T$ and $P$}{Why the mean molecular weight depends on T and P}: ionisation and molecules}
For neutral solar gas (about 91\,\% hydrogen and 9\,\% helium atoms by number, plus a small fraction of heavier elements)
$\bar\mu\approx1.3$. When a hydrogen atom is ionised it becomes two free particles, a proton and an electron, of almost
the same total mass; in fully ionised gas $\bar\mu\approx0.6$. In the coolest layers hydrogen partly forms H$_2$ molecules,
which raises $\bar\mu$. The degree of ionisation of an element in stage $i$ is given by the \textbf{Saha equation}
\begin{equation}
\frac{n_{i+1}\,n_e}{n_i}=\frac{2U_{i+1}(T)}{U_i(T)}\left(\frac{2\pi m_ek_BT}{h^2}\right)^{3/2}e^{-E_i/k_BT},\label{eq:saha}
\end{equation}
with $E_i$ the ionisation energy (13.6\,eV for hydrogen), $U_i$ the partition functions, $m_e$ the electron mass and $h$
Planck's constant. Two properties matter for us. (i) Because of the exponential, ionisation switches on over a narrow
range of temperature. (ii) The electron density $n_e$ appears on the left: at higher density there are more electrons to
recombine with, and a higher temperature is needed for the same ionisation. Hence $\bar\mu=\bar\mu(T,P)$, and
Eq.~\eqref{eq:rho_ideal} is no longer a simple proportionality. The Saha equations of all elements are coupled through the
common $n_e$ and through charge conservation ($n_e$ equals the total charge of the ions), and the molecular equilibria
(H$_2$, H$_2^+$, H$^-$) are solved together with them. This assumes \textbf{local thermodynamic equilibrium} (LTE): the
level populations and ionisation stages are those of a gas in equilibrium at the local $T$.

\subsection{\texorpdfstring{The 500\,nm extinction}{The 500 nm extinction}}
Light of wavelength 500\,nm travelling a distance $ds$ loses the fraction $\chi_{500}\,ds$ of its intensity, by absorption
and by scattering. In the solar photosphere $\chi_{500}$ is dominated by the negative hydrogen ion H$^-$ (a hydrogen atom
with a second, weakly bound electron, binding energy 0.75\,eV), whose abundance is proportional to the electron density.
The electrons come mostly from easily ionised metals (Na, Mg, Si, Ca, Fe, \dots), so $\chi_{500}$ grows steeply with
temperature and with density. The code adds the continuum contributions of H (bound--free and free--free), H$^-$,
H$_2^+$, He, He$^-$, the metals C, N, Mg, Si, Ca and Fe, Rayleigh scattering on H, He and H$_2$, and Thomson scattering
on free electrons, each counted once. Spectral lines are not included: $\chi_{500}$ is a continuum extinction.

\subsection{The table and its interpolation}
Solving the coupled Saha, molecular and charge equations at every grid point at every step would be far too slow. They
are therefore solved once, with Wittmann's equation-of-state code for solar composition (with a dedicated hot branch above
a switch temperature $T_{\rm sw}$, where the original solver fails in the hot, low-pressure corner), on a grid of
$X=\ln T$ and $Y=\ln P$, and stored in a table. For every node the table holds the three functions
\[
\ln\rho(X,Y),\qquad \ln n_e(X,Y),\qquad \ln\chi_{500}(X,Y),
\]
together with their derivatives $\partial_X$, $\partial_Y$ and $\partial_X\partial_Y$. Between the nodes each function is
evaluated by \textbf{bicubic Hermite interpolation}: a cubic polynomial in each direction that matches the stored values
and derivatives at the four corners of the cell. The result is continuous with continuous first derivatives, which is
important for the Newton-type solvers of the magnetohydrostatic modes, and it provides the derivative
$\partial\ln\rho/\partial\ln P|_T$ that they need. The table used by all modes is \code{cache/eos\_table\_fullz.h5}.
Its range of $T$ and $P$ (its \emph{envelope}) covers the input atmospheres; a point outside the envelope is reported as
an equation-of-state failure, never extrapolated.

Two derived quantities appear repeatedly:
\begin{equation}
q\equiv\frac{\rho}{P}\quad[{\rm s^2\,m^{-2}}],\qquad H\equiv\frac{P}{\rho\,g}=\frac{1}{g\,q}=\frac{k_BT}{\bar\mu\,m_{\rm u}\,g}\quad[{\rm m}].
\label{eq:H}
\end{equation}
$H$ is the \textbf{pressure scale height}. For $T=6000$\,K and $\bar\mu=1.3$, $H\approx140$\,km.
